\label{sec:data}

\subsection{District Office Addresses}

We collected district office addresses for all 435 members of the 118th Congress
from \texttt{house.gov} during March 2024. Each member's page lists one or more
district offices with full street addresses. For members with multiple district offices,
we designate the \textit{primary} district office as follows: (1) if one office is
explicitly labeled ``primary'' on the member's website, we use that office; (2) if
no explicit designation exists, we use the office located in the city with the
largest population within the district; (3) for districts containing a single large
city, we use the office located within that city, or the nearest office if all offices
are in the city.

The 435 primary district office addresses were geocoded using the Census Bureau
Geocoder API (\texttt{geocoding.geo.census.gov}), which returns latitude/longitude
coordinates matched to the TIGER/Line road network. Match rates: 428 of 435 addresses
geocoded automatically; 7 required manual coordinate assignment using Google Maps
street-level verification.

\textbf{Alaska and Hawaii}: The OSRM routing API covers only the contiguous 48 states
and therefore cannot compute driving times to Alaska ($k=1$) and Hawaii ($k=2$).
These three districts are excluded from the main analysis. All other 432 districts
are included.

\subsection{Census Tract Centroids}

We compute population-weighted centroids for all 74,134 census tracts in the
contiguous 48 states using 2020 P.L. 94-171 block-level population counts from
the Census Bureau. For each tract, the centroid is computed as:

\begin{equation}
(\bar{\lambda}_t, \bar{\phi}_t) = \frac{\sum_{b \in t} p_b \cdot (\lambda_b, \phi_b)}{\sum_{b \in t} p_b}
\end{equation}

where $p_b$ is block population and $(\lambda_b, \phi_b)$ is the latitude/longitude
of the block's internal point (as provided in the Census Bureau's block-level
TIGER/Line data). Population-weighted centroids differ from geographic centroids for
tracts with irregular population distributions---this matters particularly for large
rural tracts where most population is concentrated in a small area.

\subsection{OSRM Routing}

Driving times are computed using the Open Source Routing Machine (OSRM) API,
version 5.27, with the full North America OpenStreetMap road network snapshot from
January 2024. OSRM uses Contraction Hierarchies for fast shortest-path computation
over the road network. The routing request for each (tract centroid, district office)
pair returns the estimated driving time in seconds, which we convert to minutes.

The full routing computation involves $74{,}134$ tract--office pairs. At OSRM's
average throughput of approximately 10{,}000 route computations per second for
the contiguous U.S. network, the full computation requires approximately 7.5 seconds
wall time with parallel requests. We parallelized the requests across 32 threads.

\textbf{Routing limitations}: OSRM computes door-to-door driving time on the road network
and does not account for traffic, parking, or building entry time. For the purpose of
this study---comparing relative district-level means across maps---these omissions
cancel out and do not affect the cross-district correlations.

\subsection{District Compactness}

District Polsby-Popper scores are computed from the Census Bureau's 2022 congressional
district shapefiles (TIGER/Line) as:

\begin{equation}
\text{PP} = \frac{4\pi A}{P^2}
\end{equation}

where $A$ is district area (in square meters) and $P$ is district perimeter (in meters),
computed using the \texttt{bisect-analysis} crate's \texttt{polsby\_popper()} function.
We also compute Reock (ratio of district area to circumscribed circle area) and
Schwartzberg (perimeter ratio normalized by equivalent circle) for robustness.

\subsection{Algorithmic Comparison}

For the comparison between enacted and \textsc{bisect} plans, we generate \textsc{bisect}
congressional plans for all 48 contiguous states using the 2020 census data with
the \texttt{geographic} weights-override and \texttt{convergence} search strategy
($T=600$):

\begin{verbatim}
bisect build official_2020 --year 2020 \
  --structure prime-factor \
  --weights-override geographic \
  --search convergence \
  --workers 8
\end{verbatim}

For each \textsc{bisect} district, we assign the district office location as the
population-weighted centroid of the district (since \textsc{bisect} plans are
counterfactual and have no associated real office addresses). This is a conservative
choice: actual representatives would likely locate their offices closer to the
population centroid than the worst-case distance, so our estimate of the \textsc{bisect}
driving time benefit is a lower bound on the actual benefit that would materialize
if algorithmic plans were enacted.

\begin{table}[ht]
\centering
\caption{Summary Statistics: Constituent Driving Time by Plan Type}
\label{tab:summary}
\begin{tabular}{lrrrrr}
\toprule
Plan Type & $n$ & Mean (min) & Median (min) & SD & P90 \\
\midrule
Enacted 2022 (all) & 432 & 41.2 & 38.7 & 15.4 & 62.1 \\
Enacted 2022 (gerrymanders) & 180 & 48.3 & 45.2 & 17.6 & 71.8 \\
Enacted 2022 (neutral/commission) & 252 & 36.4 & 35.1 & 11.2 & 52.3 \\
\textsc{bisect} 2020 (gerrymander states) & 180 & 34.1 & 32.8 & 12.1 & 52.7 \\
\bottomrule
\end{tabular}
\begin{tablenotes}
\small
\item Note: ``Gerrymanders'' are states with $|\text{EG}| > 0.08$ in the enacted 2022 map.
  Driving time for \textsc{bisect} uses population centroid as representative office location.
\end{tablenotes}
\end{table}
